\section{Related work}
\label{sec:related}

\iffalse
\vspace{-1.ex}
\stitle{Optimizations on prefill \& decode disaggregation}
Both~\cite{strati2024djvu} and~\cite{zhong2024distserve} proposed prefill phase and decode phase disaggregation to XXXX.
Since these work still reserve a full copy of model parameters for disaggregated phases,
they are orthogonal to our method.
\fi

\nospacestitle{Optimizing model serving without autoscaling. }
Serving models at scale is non-trivial.
A significant body of research focuses on how to efficiently utilize GPUs to accelerate model serving~\cite{DBLP:conf/osdi/ZhongLCHZL0024,DBLP:journals/corr/abs-2404-09526,vllm,alpaserve,DBLP:conf/osdi/YuJKKC22,DBLP:conf/isca/PatelCZSGMB24,298501,kvcache},
e.g., 
Orca~\cite{DBLP:conf/osdi/YuJKKC22} proposes iterative-scheduling and selective batching. 
AlphaServe~\cite{alpaserve} employs pipeline parallelism to better handle load spikes,
but it cannot adjust pipeline instances dynamically.
%DistServe~\cite{DBLP:conf/osdi/ZhongLCHZL0024} and Splitwise~\cite{DBLP:conf/isca/PatelCZSGMB24}
%disaggregate the instances into prefill and decode pools to better
%handle the unique characteristics of these two phases.
These systems assume running on a fixed pool of GPUs,
and we have shown the
necessity of dynamically adjusting pool size and how to achieve so efficiently with {\sys}.
{\sys} complements these single-instance model serving systems with fast autoscaling mechanisms:
we build upon them for fast model serving on a single instance,
and additionally provide ultra-fast scaling when the system needs to change
the number of serving instances.

\stitle{Dynamic scaling serving instances. }
Dynamically scaling serving instances is challenging, mainly because the size of model weights is huge
and still increasing,
so loading them to the accelerators (data plane) is time-consuming.
Some existing works accelerate the loading~\cite{pipeswitch, DBLP:conf/eurosys/JeongBA23, DBLP:conf/asplos/MiaoSDXL0J24, DBLP:conf/osdi/SunHZXZL024}:
For example, both PipeSwitch~\cite{pipeswitch} and DeepPlan~\cite{DBLP:conf/eurosys/JeongBA23} leverage the layer-by-layer character of models
to overlap the inference and parameter loading to hide the loading cost. 
They only focus on host-to-device loading and such overlap is not live especially when the models are large.
SpotServe~\cite{DBLP:conf/asplos/MiaoSDXL0J24} and Llumnix~\cite{DBLP:conf/osdi/SunHZXZL024}
realize live migration but migration cannot fully unleash the computing
capabilities of both instances.
{\sys} provides a new mechanism to scale serving instances lively during parameter loading,
resulting in throughput increase even with unfinished loading,
which we have shown critical in reducing latencies under bursty workloads.

A concurrent work $\lambda$Scale~\cite{yu2025lambdascale} also focuses on using network to accelerate model autoscaling. 
The key difference is that $\lambda$Scale scatters the parameters scaling on multiple instances
to reduce the time for scaling at the cost of decreased serving throughput,
while {\sys} seamlessly scales full parameters on all instances with a similar speed,
yet does not sacrifice the throughput thanks to our multicast-chain-based scaling.
Moreover, during scaling {\sys} has a gradually increasing throughput thanks to our live scaling
while $\lambda$Scale is still a stop-the-world approach.
\iffalse
The key difference is that $\lambda$Scale mixes multiple new instances during scaling in a pipeline parallelism manner, allowing 
each instance to start serving when receiving fewer parameters (i.e., several layers of one PP stage), while {\sys} preserves an independent abstraction of each scaling instance. 
In detail, $\lambda$Scale uses multi-step binomial multicast for different pipeline stages, so multicast time increases as the number of scaling instances grows. 
Meanwhile, {\sys} employs multiple broadcast chains and leverages aggregated network bandwidth, achieving reduced multicast time even as the number of scaling instances grows. 
Moreover, $\lambda$Scale adopts a new-new collaborative execution pattern, resulting in a stepwise throughput increase 
because each new instance must transfer one PP stage before execution.
In contrast, {\sys}'s collaborative execution follows a new-old fashion, with throughput increasing gradually during the scaling process, 
further eliminates stop-of-the-world time in serverless LLM inference. 
The idea of {\sys} can be seamlessly combined with $\lambda$Scale and achieve ultra-fast autoscaling.
\fi


\stitle{Accelerating coldstart in serverless computing. }
Accelerating model scaling builds upon coldstart acceleration in
serverless computing~\cite{sock,sand,FAASM,DBLP:conf/asplos/DuYXZYQWC20,DBLP:conf/asplos/Ustiugov0KBG21,DBLP:conf/eurosys/SaxenaJSKA22,DBLP:conf/usenix/WangCTWYLDC21},
which focuses on starting general-purpose computing instances like containers. 
We built upon these works, e.g., for accelerating container startup time,
yet designed efficient network-based live autoscale tailored for model scaling
with the domain-specific knowledge of model serving.
\iffalse
For example, serving instances still need to start up containers first.
A key difference is that current general-purpose acceleration techniques,
e.g., on-demand loading~\cite{DBLP:conf/eurosys/WangHW19,mitosis}
that skips loading all data, are not applicable to model serving
because models require all parameters for execution.
{\sys} overcomes this challenge by leveraging domain-specific knowledge of the layer-by-layer execution
pattern of models for acceleration.
Meanwhile, our multicast is tailored for model serving
as the instance sizes are orders of magnitude larger than the container sizes in traditional serverless computing.
\fi